\chapter*{Zakończenie}
\addcontentsline{toc}{chapter}{Zakończenie}
Celem pracy było utworzenie aplikacji pośredniczącej między zewnętrznymi systemami, a Krajowym Systemem e-Faktur, ze szczególnym naciskiem na prostotę wdrożenia oraz obsługi. 

Aplikacja została utworzona z wykorzystaniem języka Go, z bazą danych SQLite oraz panelem użytkownika wykorzystującym HTMX. Wykorzystane technologie pozwoliły na spełnienie większości założeń. Cel pracy został osiągnięty, jak zostało zademonstrowane na końcu rozdziału trzeciego, w dwóch scenariuszach: jednym prezentującym przykład poprawnie przetworzonej faktury i jednym z fakturą odrzuconą przez KSeF. W wyniku pracy powstał pojedynczy plik wykonywalny.

Możliwe byłoby rozszerzenie funkcjonalności aplikacji, przede wszystkim poprzez wyeliminowanie ograniczeń wyliczonych na końcu rozdziału trzeciego: wsparcie dla trybów offline oraz offline24, wsparcie dla wielu podmiotów, czy poszerzenie pokrycia schematu FA(3).

Pełne repozytorium z aplikacją można znaleźć pod adresem \url{https://github.com/brunonsocha/ksef-integration}.
